\documentclass[12pt,a4paper]{ctexart}

\usepackage[margin=1.65cm,top=1.45cm,bottom=1.45cm]{geometry}
\usepackage{amsmath,amssymb}
\usepackage{tikz}
\usepackage{enumitem}

\pagestyle{empty}

\setlength{\parindent}{0pt}
\setlength{\parskip}{0.35em}
\setlist[enumerate]{leftmargin=1.8em,itemsep=0.45em,topsep=0.2em}

\newcommand{\lessonTitle}[2]{
  \begin{center}
    {\LARGE\bfseries #1}
  \end{center}
  \vspace{0.2em}
  \hrule
  \vspace{0.8em}
}

\newcommand{\questionBox}[2]{
  \textbf{\large #1}

  \vspace{0.25em}
  #2
  \vspace{0.45em}
  \hrule
}

\newcommand{\writeArea}[1]{
  \vspace{0.75em}
  \vspace*{#1cm}
}

\begin{document}

\lessonTitle{函数与导数课堂讲义}{}

\textbf{\large 知识点梳理一：函数的三要素与基本表示}

\begin{enumerate}
  \item \textbf{函数的本质：}函数描述的是“一个输入对应一个唯一输出”的规则。若 $x$ 在集合 $A$ 中取值，按确定规则 $f$ 在集合 $B$ 中得到唯一的 $y$，就记作 $y=f(x)$。
  \item \textbf{三要素：}函数由定义域、对应法则、值域共同决定。做函数题先看 $x$ 能不能取，再看 $x$ 取定后 $y$ 怎样变化，最后看 $y$ 能取到哪些值。
  \item \textbf{定义域：}允许代入的全部 $x$ 的集合。常见限制包括分母不为 $0$、偶次根式被开方数非负、对数真数为正、实际问题中的长度/时间/数量限制等。
  \item \textbf{对应法则：}把 $x$ 变成 $y$ 的规则，可以由解析式、图像、表格或文字给出。两个函数相同，必须定义域相同且对应法则相同。
  \item \textbf{值域：}函数值 $y$ 能达到的全部结果。求值域常用图像观察、单调性、最值、不等式、换元后限制新变量范围等方法。
  \item \textbf{表示方式：}解析式适合计算，图像适合看变化趋势，表格适合看离散数据，文字规则适合实际情境。不同表示之间可以互相翻译。
\end{enumerate}

\vspace{0.5em}
\hrule

\newpage
\lessonTitle{函数与导数课堂讲义}{}

\textbf{\large 知识点梳理二：函数图像的读法}

\begin{enumerate}
  \item \textbf{图像上的点：}点 $(x_0,y_0)$ 在函数图像上，表示 $y_0=f(x_0)$。横坐标看输入，纵坐标看输出。
  \item \textbf{从图像读定义域和值域：}图像在水平方向覆盖的 $x$ 范围是定义域，在竖直方向覆盖的 $y$ 范围是值域。空心点、断点、端点是否取到要特别看清。
  \item \textbf{单调性：}从左往右看，图像上升表示函数递增，图像下降表示函数递减。单调区间必须写在定义域内部。
  \item \textbf{最值：}最高点对应最大值，最低点对应最小值；若只在某个区间内研究，还要检查区间端点。
  \item \textbf{零点：}图像与 $x$ 轴交点的横坐标就是零点，也就是方程 $f(x)=0$ 的实根。
  \item \textbf{对称与周期：}关于 $y$ 轴对称常对应偶函数，关于原点对称常对应奇函数，图像重复出现时可能存在周期。
\end{enumerate}

\begin{center}
\begin{tikzpicture}[scale=0.82]
  \draw[->] (-3.7,0) -- (3.7,0) node[right] {$x$};
  \draw[->] (0,-1.8) -- (0,2.6) node[above] {$y$};
  \draw[line width=0.9pt,domain=-3.0:3.0,samples=120] plot (\x,{0.12*(\x+2.3)*(\x+0.3)*(\x-2.1)+0.15});
  \fill (-2.3,0) circle (1.3pt) node[below] {零点};
  \fill (-0.3,0) circle (1.3pt) node[above] {零点};
  \fill (2.1,0) circle (1.3pt) node[below] {零点};
  \fill (-1.55,1.12) circle (1.3pt) node[above] {局部最高};
  \fill (1.25,-0.95) circle (1.3pt) node[below] {局部最低};
  \draw[dashed] (-3.0,-1.5) -- (-3.0,2.2) node[above] {端点};
  \draw[dashed] (3.0,-1.5) -- (3.0,2.2) node[above] {端点};
\end{tikzpicture}
\end{center}

\hrule

\newpage
\lessonTitle{函数与导数课堂讲义}{}

\textbf{\large 知识点梳理三：二次函数的图像与性质}

\begin{enumerate}
  \item \textbf{三种常见形式：}一般式 $y=ax^2+bx+c\ (a\ne0)$；顶点式 $y=a(x-h)^2+k$；交点式 $y=a(x-x_1)(x-x_2)$。形式不同，读取信息的重点不同。
  \item \textbf{开口方向与宽窄：}$a>0$ 开口向上，$a<0$ 开口向下；$|a|$ 越大，图像越窄，$|a|$ 越小，图像越宽。
  \item \textbf{对称轴：}$x=-\dfrac{b}{2a}$，或由顶点式直接读 $x=h$。对称轴左右等距离的点函数值相等。
  \item \textbf{顶点：}一般式顶点为 $\left(-\dfrac{b}{2a},\dfrac{4ac-b^2}{4a}\right)$，顶点式顶点为 $(h,k)$。顶点是整条抛物线的最高点或最低点。
  \item \textbf{判别式与根：}$\Delta=b^2-4ac$。$\Delta>0$ 有两个不同实根，$\Delta=0$ 有一个重根并与 $x$ 轴相切，$\Delta<0$ 没有实根。
  \item \textbf{单调区间与区间最值：}开口向上时，对称轴左侧递减、右侧递增；开口向下相反。在有限区间上求最值时，要比较端点和对称轴位置。
\end{enumerate}

\begin{center}
\begin{tikzpicture}[scale=0.86]
  \draw[->] (-3.2,0) -- (3.2,0) node[right] {$x$};
  \draw[->] (0,-1.0) -- (0,3.5) node[above] {$y$};
  \draw[dashed] (0.8,-0.75) -- (0.8,3.1) node[above] {对称轴};
  \draw[line width=0.9pt,domain=-1.3:2.9,samples=100] plot (\x,{0.55*(\x-0.8)^2-0.55});
  \fill (0.8,-0.55) circle (1.5pt) node[below right] {顶点};
  \fill (-0.2,0) circle (1.2pt) node[above left] {交点};
  \fill (1.8,0) circle (1.2pt) node[above right] {交点};
  \node at (-1.55,2.45) {$a>0$，开口向上};
\end{tikzpicture}
\end{center}

\hrule

\newpage
\lessonTitle{函数与导数课堂讲义}{}

\textbf{\large 知识点梳理四：最值与极值的关系}

\begin{enumerate}
  \item \textbf{最大值与最小值：}是在指定范围内函数值能达到的最大或最小结果，是“全局”概念。没有指定范围时，必须先确认研究的是整个定义域还是某个区间。
  \item \textbf{极大值与极小值：}只比较某点附近左右邻近位置的函数值，是“局部”概念。极大值不一定是最大值，极小值也不一定是最小值。
  \item \textbf{区间最值候选点：}闭区间上的最值通常从端点、区间内部极值点、不可导点、分段连接点、边界临界点中产生。
  \item \textbf{端点不能漏：}导数只反映区间内部的变化，端点不需要满足 $f'(x)=0$，但端点函数值完全可能成为最大值或最小值。
  \item \textbf{开区间与闭区间：}闭区间连续函数一定能取到最大值和最小值；开区间或无穷区间可能只有上确界/下确界而取不到最值。
  \item \textbf{常见判断顺序：}先定定义域和研究区间，再找候选点，计算候选点函数值，最后比较大小并确认是否取到。
\end{enumerate}

\vspace{0.5em}
\hrule

\newpage
\lessonTitle{函数与导数课堂讲义}{}

\textbf{\large 知识点梳理五：极值的必要条件与充分条件}

\begin{enumerate}
  \item \textbf{必要条件：}若函数在可导点 $x_0$ 处取得极值，通常有 $f'(x_0)=0$。这说明“有极值 $\Rightarrow$ 导数可能为零”，不是反过来。
  \item \textbf{必要条件不充分：}$f'(x_0)=0$ 只能说明该点可能是极值点，还要看左右变化。例如导数左右不变号时，函数可能只是“平着穿过”。
  \item \textbf{一阶导充分条件：}若 $f'(x)$ 在 $x_0$ 左右由正变负，则 $x_0$ 处为极大值；由负变正，则为极小值。
  \item \textbf{单调性判定：}也可以不直接背结论，而是看 $x_0$ 左边函数增减、右边函数增减：先增后减是极大，先减后增是极小。
  \item \textbf{二阶导辅助：}在允许使用二阶导时，若 $f'(x_0)=0$ 且 $f''(x_0)<0$，可能为极大值；若 $f''(x_0)>0$，可能为极小值。若 $f''(x_0)=0$，还需另判。
  \item \textbf{易错点：}不要把“极值点的横坐标”“极值”“最值”混成同一个量；也不要只找到 $f'(x)=0$ 就直接下结论。
\end{enumerate}

\begin{center}
\begin{tikzpicture}[scale=0.82]
  \draw[->] (-3.4,0) -- (3.4,0) node[right] {$x$};
  \draw[->] (0,-1.6) -- (0,2.4) node[above] {$y$};
  \draw[line width=0.9pt,domain=-2.8:2.8,samples=120] plot (\x,{0.18*(\x+1.6)*(\x-0.2)*(\x-2.2)+0.45});
  \fill (-1.45,1.37) circle (1.4pt) node[above] {极大值};
  \fill (1.55,-0.85) circle (1.4pt) node[below] {极小值};
  \draw[->] (-2.5,0.7) -- (-1.75,1.18) node[midway,above] {增};
  \draw[->] (-1.1,1.18) -- (-0.35,0.63) node[midway,above] {减};
  \draw[->] (0.6,-0.35) -- (1.25,-0.75) node[midway,below] {减};
  \draw[->] (1.85,-0.75) -- (2.55,0.25) node[midway,right] {增};
\end{tikzpicture}
\end{center}

\hrule

\newpage
\lessonTitle{函数与导数课堂讲义}{}

\textbf{\large 知识点梳理六：零点、方程根与图像交点}

\begin{enumerate}
  \item \textbf{零点定义：}若 $f(x_0)=0$，则 $x_0$ 是函数 $f(x)$ 的零点。零点是横坐标，不是图像上的整个点。
  \item \textbf{零点与方程：}函数 $f(x)$ 的零点等价于方程 $f(x)=0$ 的实根，也等价于图像 $y=f(x)$ 与 $x$ 轴交点的横坐标。
  \item \textbf{存在性：}若函数在 $[a,b]$ 上连续，且 $f(a)f(b)<0$，则 $(a,b)$ 内至少有一个零点。这个结论保证存在，不直接保证唯一。
  \item \textbf{唯一性：}若再结合函数在区间上严格单调，就能把“至少一个”推进到“恰有一个”。
  \item \textbf{图像状态：}穿过 $x$ 轴时两侧函数值变号；相切时可能不变号；整段在轴上方或下方时没有零点。
  \item \textbf{含参临界：}参数改变图像位置、斜率或形状时，零点个数的变化常发生在相切、过端点、极值点落在 $x$ 轴上、定义域边界变化等状态。
\end{enumerate}

\begin{center}
\begin{tikzpicture}[scale=0.78]
  \draw[->] (-3.2,0) -- (3.2,0) node[right] {$x$};
  \draw[->] (0,-1.35) -- (0,2.2) node[above] {$y$};
  \draw[line width=0.85pt,domain=-2.55:2.45,samples=120] plot (\x,{0.22*(\x+2.0)*(\x-0.2)*(\x-1.7)});
  \fill (-2,0) circle (1.3pt) node[below] {穿过};
  \fill (-0.2,0) circle (1.3pt) node[above] {穿过};
  \fill (1.7,0) circle (1.3pt) node[below] {穿过};
  \begin{scope}[xshift=8.0cm]
    \draw[->] (-2.0,0) -- (2.0,0) node[right] {$x$};
    \draw[->] (0,-0.7) -- (0,2.2) node[above] {$y$};
    \draw[line width=0.85pt,domain=-1.55:1.55,samples=100] plot (\x,{0.65*\x*\x});
    \fill (0,0) circle (1.3pt) node[below right] {相切};
  \end{scope}
\end{tikzpicture}
\end{center}

\hrule

\newpage
\lessonTitle{函数与导数课堂讲义}{}

\questionBox{第1题：概念与基础计算}{
已知函数 $f(x)=x^2-2ax+1$。
\begin{enumerate}
  \item 写出函数图象的对称轴。
  \item 当 $x\in[0,3]$ 时，讨论 $f(x)$ 的最小值与参数 $a$ 的关系。
\end{enumerate}
}

\writeArea{15.4}

\newpage
\lessonTitle{函数与导数课堂讲义}{}

\questionBox{第2题：函数零点与参数}{
设函数 $g(x)=\ln x-kx+1$，其中 $k$ 为实数。
\begin{enumerate}
  \item 当 $k=1$ 时，判断 $g(x)$ 在 $(0,+\infty)$ 上的零点个数。
  \item 若函数 $g(x)$ 在 $(0,+\infty)$ 上恰有两个零点，求 $k$ 的取值范围。
\end{enumerate}
}

\vspace{0.5em}
\begin{center}
\begin{tikzpicture}[scale=0.95]
  \draw[->,line width=0.65pt] (-0.2,0) -- (7.2,0) node[right] {$x$};
  \draw[->,line width=0.65pt] (0,-2.0) -- (0,3.2) node[above] {$y$};
  \draw[line width=0.95pt,domain=0.25:6.6,samples=120] plot (\x,{ln(\x)+1-0.55*\x});
  \draw[gray!55,line width=0.55pt] (1,0.08) -- (1,-0.08) node[below=4pt] {$1$};
  \node at (4.95,1.2) {$y=\ln x-kx+1$};
\end{tikzpicture}
\end{center}

\writeArea{10.7}

\newpage
\lessonTitle{函数与导数课堂讲义}{}

\questionBox{第3题：同类小题组}{
\begin{enumerate}
  \item 已知 $\vec a=(2,-1)$，$\vec b=(m,3)$，若 $\vec a\perp \vec b$，求 $m$。
  \item 已知复数 $z=1+2i$，求 $z\overline z$。
  \item 正方体 $ABCD-A_1B_1C_1D_1$ 中，判断直线 $AC$ 与平面 $BDD_1B_1$ 的位置关系。
\end{enumerate}
}

\writeArea{15.3}

\newpage
\lessonTitle{函数与导数课堂讲义}{}

\textbf{\large 常用套路模板总结}

\vspace{0.3em}
\begin{enumerate}
  \item \textbf{区间二次函数最值：}适用于开口方向已知、区间固定、参数影响对称轴的题。先写对称轴，再比较它与区间左端点、右端点的位置。
  \item \textbf{零点个数讨论：}适用于函数连续且可通过单调性或图像趋势判断交点的题。先确定定义域，再看导数或图像变化，最后结合端点/极值判断零点个数。
  \item \textbf{容易失效的情况：}定义域变化、参数同时影响区间和函数式、图像不连续时，不能直接套上述流程。
\end{enumerate}

\end{document}
